\paragraph{E.4l (Proved): the actual twist-zero support and every finite-prime open.}
The obstruction in E.4j has an explicit replacement. Throughout the next
five paragraphs put
\[
 X=\operatorname{Spec}\mathbb Z,\qquad
 U_S=\operatorname{Spec}\mathbb Z[S^{-1}],\qquad
 \overline U_S=\overline X-S,
\]
where $S$ is a finite set of rational primes; thus the real place is
retained. On $\overline U_S$ use the restriction of the Artin--Verdier
complex $\mathbb Z(n)_{\overline X}$ defined in
\citet[\S6.2 and Definition 6.5]{D_FlachMorin2018}, and write
\[
 A_S^n=R\Gamma(\overline U_{S,\mathrm{et}},
                     \mathbb Z(n)_{\overline X}|_{\overline U_S}).
\]
This notation specifies a complex, including its real-place convention.
It is not a claim that the manuscript defined these coefficients.
By Proposition 6.11 of that paper, $A_S^0$ is the cohomology of the
constant sheaf $\mathbb Z$. Define the discrete torsion group
\[
 D_p=\operatorname{Hom}_{\mathrm{cts}}
                   (\mathbb Z_p^\times,\mathbb Q/\mathbb Z).
\]
There is a local support isomorphism
\[
 R\Gamma_{\{p\}}(X_{\mathrm{et}},\mathbb Z)\simeq D_p[-3].
 \tag{E.4.13}
\]
Its tame subgroup is
$\operatorname{Hom}(\mathbb F_p^\times,\mathbb Q/\mathbb Z)$, induced by
reduction of units. The quotient is
$\operatorname{Hom}_{\mathrm{cts}}(1+p\mathbb Z_p,\mathbb Q/\mathbb Z)$.
For odd $p$ this quotient is $\mathbb Q_p/\mathbb Z_p$; for $p=2$ it is
$\mathbb Z/2\oplus\mathbb Q_2/\mathbb Z_2$.
Consequently, as abstract groups,
\[
 D_p\simeq
 \begin{cases}
  \mathbb Z/(p-1)\oplus\mathbb Q_p/\mathbb Z_p,&p\ne2,\\
  \mathbb Z/2\oplus\mathbb Q_2/\mathbb Z_2,&p=2.
 \end{cases}
 \tag{E.4.14}
\]
Only the finite tame summand is detected by
$C_p^{-1}[-2]$ from E.4a. The remaining summand is actual wild support.

For nonempty $S$, set
\[
 D_S^+=\operatorname{Hom}_{\mathrm{cts}}
 \left(
   \left(\prod_{p\in S}\mathbb Z_p^\times\right)/
                  \langle(-1)_{p\in S}\rangle,\ \mathbb Q/\mathbb Z
 \right).
\]
Then every cohomology group is given by
\[
 H^i(A_S^0)=
 \begin{cases}
  \mathbb Z,&i=0,\\
  D_S^+,&i=2,\\
  0,&\text{otherwise}.
 \end{cases}
 \qquad(S\ne\varnothing).
 \tag{E.4.15}
\]
For $S=\varnothing$, $A_\varnothing^0\simeq C_{\overline X}^0$
has the groups E.4b. In particular,
\[
 A_{\{2\}}^0\simeq
      \mathbb Z\oplus(\mathbb Q_2/\mathbb Z_2)[-2].
 \tag{E.4.16}
\]

\begin{proof}
Etale excision reduces (E.4.13) to the excellent henselian discrete
valuation ring $R=\mathbb Z_{(p)}^h$. The local duality pairing
of \citet[II, Theorem 1.8, with its pairing defined on p.~152]{D_MilneADT2006}
identifies $H^3_{\{p\}}(\operatorname{Spec}R,\mathbb Z)$ with the
continuous dual of the completion of $R^\times$, hence with $D_p$.
Here $\operatorname{Ext}^r_R(\mathbb Z,\mathbb G_m)
=H^r(R,\mathbb G_m)=0$ for $r>0$ by II, Proposition 1.5(a).
For completeness, the vanishing in all degrees and the shift can also
be read directly from localization. II, Corollary 1.2 gives
$H^0(R,\mathbb Z)=\mathbb Z$ and $H^2(R,\mathbb Z)=\mathbb Q/\mathbb Z$,
with the other groups zero. If $K=\operatorname{Frac}R$, local class
field theory and the sequence
$0\to\mathbb Z\to\mathbb Q\to\mathbb Q/\mathbb Z\to0$ give
\[
 H^2(K,\mathbb Z)
   =\operatorname{Hom}_{\mathrm{cts}}(\widehat{K^\times},
                                      \mathbb Q/\mathbb Z).
\]
The degree-two restriction from $R$ is the unramified character
subgroup, dual to the valuation $K^\times\to\mathbb Z$.
Also $H^3(K,\mathbb Z)=H^2(K,\mathbb Q/\mathbb Z)=0$:
finite local duality identifies the latter group's dual with the
inverse limit of finite groups of roots of unity in $K$, which is
zero because $K$ has only finitely many roots of unity.
Higher torsion cohomology vanishes by local cohomological dimension
two. The remaining quotient is precisely the dual of the unit
group. This proves (E.4.13), including that there are no omitted
cohomology groups. The equality of the henselian and completed
Galois groups is recalled in
\citet[I, proof of Theorem 2.14, pp.~36--37]{D_MilneADT2006}.

For odd $p$, the Teichmuller decomposition
$\mathbb Z_p^\times=\mu_{p-1}\times(1+p\mathbb Z_p)$ and the
$p$-adic logarithm identify the second factor with $\mathbb Z_p$.
For $p=2$, use
$\mathbb Z_2^\times=\{\pm1\}\times(1+4\mathbb Z_2)$ and the
logarithm on the second factor. This proves (E.4.14) and the
assertions about the tame subgroup.

The rational correction in Definition 3.6 of
\citet{D_FlachMorin2018} is zero in twist zero:
$R\Gamma(X,\mathbb Q(1))=0$, since $\mathbb Z(1)=\mathbb G_m[-1]$,
$\mathbb Z^\times$ and $\operatorname{Pic}(\mathbb Z)$ have zero
rationalization, and the higher $\mathbb G_m$ groups are torsion
by \citet[II, Proposition 2.1]{D_MilneADT2006}.
Thus $A_\varnothing^0=C_{\overline X}^0$.
Localization and E.4b now give, for any finite $S$,
\[
 0\longrightarrow H^2(A_S^0)\longrightarrow
       \bigoplus_{p\in S}D_p
   \xrightarrow{\partial_S}\mathbb Z/2
   \longrightarrow H^3(A_S^0)\longrightarrow0,
 \tag{E.4.17}
\]
as well as $H^0(A_S^0)=\mathbb Z$ and vanishing in all other degrees.

We identify the kernel and show surjectivity rather than assume
the boundary map. Constant rational coefficients have no positive
cohomology on $\overline U_S$: on the normal etale charts they are
the direct image of the constant rational sheaf at the generic
point, whose positive Galois cohomology vanishes by averaging over
finite quotients; the real stabilizer has the same vanishing.
It follows that
$H^2(A_S^0)=H^1(\overline U_{S,\mathrm{et}},\mathbb Q/\mathbb Z)$.
Finite abelian torsors on this Artin--Verdier site are the abelian
extensions of $\mathbb Q$ unramified outside $S$ that are totally
real. By the Kronecker--Weber theorem, with the ramification
control in its proof
\citep[I, Theorem 4.16 and Lemma 4.17, pp.~49--50]{D_MilneCFT2020},
their maximal Galois group is
\[
 \left(\prod_{p\in S}\mathbb Z_p^\times\right)/
                      \langle(-1)_{p\in S}\rangle.
\]
Indeed the maximal cyclotomic extension unramified outside $S$
has the product before quotienting, and complex conjugation is
the indicated diagonal element. Naturality of the coefficient
and localization sequences identifies the injection in (E.4.17)
with restriction of these characters to the local unit groups.
Its image is the annihilator of the diagonal element. The
quotient therefore has order two when $S\ne\varnothing$.
Consequently $\partial_S$ is surjective and is, under these
identifications, evaluation on that element. This proves
(E.4.15). For $S=\{2\}$ the quotient of the unit group by
$\{\pm1\}$ is $1+4\mathbb Z_2\simeq\mathbb Z_2$.
Finally the two nonzero cohomology groups in (E.4.16) split
uniquely in $D(\mathbb Z)$: the relevant ambiguity and obstruction
are $\operatorname{Ext}^2$ and $\operatorname{Ext}^3$, both zero.
\end{proof}

\paragraph{E.4m (Proved): the exact extra local term in every negative twist.}
For $n<0$, the coefficient convention of
\citet[\S3.1 and \S6.2]{D_FlachMorin2018} on $X_{\mathrm{et}}$ is
\[
 \mathbb Z(n)=
   \bigoplus_\ell j_{\ell,!}(\mathbb Q_\ell/\mathbb Z_\ell(n))[-1],
 \qquad j_\ell:X[1/\ell]\hookrightarrow X.
 \tag{E.4.18}
\]
With this convention there is a primary decomposition of the
actual support complex
\[
 R\Gamma_{\{p\}}(X_{\mathrm{et}},\mathbb Z(n))
 \simeq
 C_p^{\,n-1}[-2]\ \oplus\
 R\Gamma(\mathbb Q_p,\mathbb Q_p/\mathbb Z_p(n))[-2],
 \qquad n<0.
 \tag{E.4.19}
\]
In particular the second summand cannot be discarded when
preserving all finite coefficients.

\begin{proof}
Work again over $R=\mathbb Z_{(p)}^h$. For $\ell\ne p$ the
cyclotomic $\ell$-primary sheaf is unramified and the codimension-one
prime-to-$p$ purity calculation gives
\[
 R\Gamma_{\{p\}}(R,\mathbb Q_\ell/\mathbb Z_\ell(n)[-1])
  \simeq
 R\Gamma(\mathbb F_p,\mathbb Q_\ell/\mathbb Z_\ell(n-1))[-3].
\]
Equivalently, its local inertia computation uses the tame
$\mathbb Z_\ell(1)$ quotient: its first cohomology is the twist
$n-1$, its higher $\ell$-primary cohomology vanishes, and the
support and coefficient shifts supply the two additional degrees.
The sum of these terms is $C_p^{n-1}[-2]$ by the
negative-twist calculation E.4a. Cohomology and support commute
with these filtered sums; the local continuity statement is also
\citet[Lemma 2.1]{D_GeisserDedekind2004}.

For the $p$-primary summand put
$j:\operatorname{Spec}K\hookrightarrow\operatorname{Spec}R$.
The exact assertion
\[
 R\Gamma(R,j_!F)=0,\qquad
 R\Gamma_{\{p\}}(R,j_!F)\simeq R\Gamma(K,F)[-1]
\]
for an etale sheaf $F$ is
\citet[II, Proposition 1.1(a)]{D_MilneADT2006}.
Apply it to $F=\mathbb Q_p/\mathbb Z_p(n)$ and include the
shift $[-1]$ already present in (E.4.18). The result is
$R\Gamma(K,\mathbb Q_p/\mathbb Z_p(n))[-2]$.
The henselian-completion equivalence of finite extensions,
and then continuity for the torsion sheaf, replace $K$ by
$\mathbb Q_p$. This proves (E.4.19).
These computations take place at a finite closed point and
therefore are unchanged by the real-place correction.
\end{proof}

\paragraph{E.4n (Proved): full integral purity for positive twists at least two on these opens.}
The failure above does not extend to every weight. If $n\ge2$,
then for the closed immersion $i_p:\operatorname{Spec}\mathbb F_p
\hookrightarrow X$ the usual cycle-class purity morphism is a
quasi-isomorphism
\[
 \mathbb Z(n-1)_{\mathbb F_p}[-2]
       \xrightarrow{\ \sim\ } Ri_p^!\mathbb Z(n)_X.
 \tag{E.4.20}
\]
Consequently its derived global support is $C_p^{n-1}[-2]$,
including the characteristic-prime part, and restriction from
$\overline U_S$ to $\overline U_T$, $S\subset T$, has the
ordinary purity triangle in these weights.

\begin{proof}
There is a direct source with precisely the needed bound:
\citet[Proposition 5.1, p.~789]{D_GeisserDedekind2004}
proves full integral purity for $n\ge\operatorname{cd}_p(U_{\mathrm{et}})$,
where $U$ is the generic fibre of an essentially smooth scheme over
a discrete valuation ring of residue characteristic $p$.
Its historical Bloch--Kato hypothesis is now the norm-residue
theorem \citep[Theorems 6.16--6.17, pp.~436--437]{D_Voevodsky2011}.
By etale excision we may use the base ring
$R=\mathbb Z_{(p)}^h$ and the identity smooth morphism.
The generic field $K$ has the same absolute Galois group as
$\mathbb Q_p$, which has cohomological dimension two for every
torsion prime, including $p$. This follows from finite local
duality and is retained by finite extensions
\citep[I, Theorems 2.1 and 2.14]{D_MilneADT2006}.
Thus Geisser's inequality holds for every $n\ge2$.
His morphism is the ordinary purity morphism, so the resulting
local quasi-isomorphisms are the stalkwise isomorphisms in
(E.4.20), not a newly chosen map.

One can check the potentially delicate $p$-primary part separately.
\citet[Corollary 4.4, p.~786]{D_GeisserDedekind2004}
verifies the vanishing hypothesis in
\citet[Lemma 7.7, pp.~1543--1545]{D_FlachMorin2018}.
For this smooth discrete valuation ring its special fibre is a
normal crossing divisor. At each exponent $p^r$ the lemma's
localization proof gives
\[
 i_{p,*}\mathbb Z/p^r(n-1)[-2]\longrightarrow
 \mathbb Z/p^r(n)\longrightarrow
 \tau_{\le n}Rj_*\mu_{p^r}^{\otimes n}\longrightarrow.
 \tag{E.4.21}
\]
The argument is made with the cone of multiplication by $p^r$;
no inference from an isomorphism of pro-objects to each level
is required here. The first term vanishes because a perfect
field of dimension zero has
$W_r\Omega_{\mathbb F_p,\log}^{n-1}=0$ for $n-1>0$.
The stalks of $R^a j_*\mu_{p^r}^{\otimes n}$ for $a>1$
vanish: the strict henselian generic field has cohomological
dimension at most one
\citep[II, Lemma 1.3 and its proof, p.~150]{D_MilneADT2006}.
The truncation in (E.4.21) is therefore unnecessary.
It follows that $Ri_p^!\mathbb Z/p^r(n)=0$, exactly as
required by (E.4.20) at the residue-characteristic prime.
This also explains why the same argument does not give
twist-zero purity.

For $n-1>0$, ordinary etale motivic and Weil-etale cohomology
of the finite point agree. Indeed the comparison triangle of
\citet[Theorem 7.1]{D_Geisser2004} has the rational splitting
written in \citet[Proposition 3.17, equation (41), p.~1459]{D_FlachMorin2018}:
the rational Weil complex is the direct sum of the rational
etale complex and its shift by $[-1]$. E.4a makes that Weil
complex rationally zero, so both summands vanish. The integral
comparison triangle then gives the claimed isomorphism.
Hence the global support of (E.4.20) is the asserted
$C_p^{n-1}[-2]$.
\end{proof}

\paragraph{E.4o (Proved): the twist-one open complex has an actual finite-coefficient comparison.}
The complex constructed in E.4e can now be compared, on these
specific opens and with the specified real convention, to the
ordinary etale motivic complex. This is a sharper result than
merely comparing the displayed groups. For nonempty $S$ put
\[
 \begin{split}
 E_S&=\mathbb Z[S^{-1}]^\times=\{\pm1\}\oplus\mathbb Z^S,\\
 B_S&=\ker\left((\mathbb Q/\mathbb Z)^S
                     \xrightarrow{\sum}\mathbb Q/\mathbb Z\right),\\
 V_S&=\ker\left(\mathbb Q^S\xrightarrow{\sum}\mathbb Q\right),
 \qquad
 \Lambda_S=\ker\left(\mathbb Z^S\xrightarrow{\sum}\mathbb Z\right).
 \end{split}
 \tag{E.4.22}
\]
The unit identification uses the actual rational primes as
positive generators. There is an exact sequence
\[
 0\longrightarrow\Lambda_S\longrightarrow V_S
       \longrightarrow B_S\longrightarrow0.
 \tag{E.4.23}
\]
Then
\[
 A_S^1\simeq E_S[-1]\oplus B_S[-3],
 \qquad
 C_S^1\simeq
 \operatorname{Cone}\left(V_S[-3]\longrightarrow A_S^1\right),
 \quad S\ne\varnothing,
 \tag{E.4.24}
\]
where the arrow is zero on $E_S[-1]$ and the quotient in
(E.4.23) on $B_S[-3]$. The second isomorphism identifies this
cone with the complex in E.4e, compatibly with every restriction
$S\subset T$. In particular, for every positive integer $m$,
\[
 C_S^1\otimes^{\mathbb L}\mathbb Z/m
   \simeq
 R\Gamma(\overline U_{S,\mathrm{et}},\mathbb Z/m(1)),
 \tag{E.4.25}
\]
with $\mathbb Z/m(1)$ the derived reduction of the stipulated
Artin--Verdier complex.

\begin{proof}
Proposition 6.11 of \citet{D_FlachMorin2018} gives
$\mathbb Z(1)=\phi_*\mathbb G_m[-1]$ on the compactified site.
The global $\mathbb G_m$ calculation of
\citet[II, Proposition 2.1 and Remark 2.2(a),
pp.~163--165]{D_MilneADT2006} identifies the ordinary Brauer
group of $U_S$ with the kernel of the sum of invariants from
the removed finite places and the real place.
The real correction removes the real invariant. Explicitly,
the triangle from Proposition 6.10 of
\citet{D_FlachMorin2018} compares $\phi_*\mathbb G_m$ with
$R\phi_*\mathbb G_m$ and the positive-degree real cohomology.
The map $\operatorname{Br}(U_S)\to\operatorname{Br}(\mathbb R)
=\mathbb Z/2$ is surjective for $S\ne\varnothing$, since a
finite removed place can carry invariant $1/2$.
Its kernel is $B_S$. Higher groups cancel with their real
counterparts, while
$H^0(\overline U_S,\phi_*\mathbb G_m)=E_S$ and
$H^1(\overline U_S,\phi_*\mathbb G_m)
=\operatorname{Pic}(U_S)=0$.
Thus $A_S^1$ has only $E_S$ in degree one and $B_S$ in
degree three. Its splitting in (E.4.24) is unique because
$\operatorname{Ext}^{a}_{\mathbb Z}(-,-)=0$ for $a\ge2$.

To verify (E.4.23), lift an element of $B_S$ to a rational
vector. Its sum is integral; subtract that integer from one
coordinate to obtain a lift in $V_S$. The kernel is exactly
$\Lambda_S$. The cone in (E.4.24) consequently has $E_S$ in
degree one and $\Lambda_S$ in degree two, with no other
cohomology. Since $\Lambda_S$ is free, its splitting is also
unique: the possible ambiguity
$\operatorname{Ext}^1(\Lambda_S,E_S)$ is zero.
Its groups, identifications and restrictions are therefore
precisely those of the integral complex in E.4e.

For the proper member the corresponding data are
\[
 A_\varnothing^1\simeq
  (\mathbb Z/2)[-1]\oplus(\mathbb Q/\mathbb Z)[-4],
 \qquad
 \mathbb Q[-4]\longrightarrow A_\varnothing^1,
 \tag{E.4.26}
\]
where the map is the quotient $\mathbb Q\to\mathbb Q/\mathbb Z$
on the second summand. These are the Artin--Verdier trace and
the defining rational correction of
\citet[Theorem 3.5, Definition 3.6 and \S5.8.3]{D_FlachMorin2018}.
Its cone is $(\mathbb Z/2)[-1]\oplus\mathbb Z[-3]$.
Restriction from the proper member to a nonempty $S$ is the
unit inclusion and zero on the other displayed summands:
$\operatorname{Ext}^1(\mathbb Q/\mathbb Z,B_S)=0$ because
$B_S$ is divisible. Likewise the map of correction complexes
$\mathbb Q[-4]\to V_S[-3]$ is zero.
For two nonempty sets use the coordinate inclusions in
(E.4.22)--(E.4.23). These give commuting diagrams of actual
complexes and hence coherent cone maps, including the empty
set. They agree with E.4e; any possible off-diagonal map
between its displayed free cohomology groups vanishes.
Finally every correction complex here is a rational vector
space in one degree, so its derived reduction modulo $m$
is zero. Tensoring the defining cone triangle proves
(E.4.25), naturally for all these restriction maps.
\end{proof}

\paragraph{E.4p (Proved with an explicitly restricted scope): a coherent all-twist integral extension with the necessary wild terms.}
There is now an explicit family of bounded integral complexes,
with restrictions for all finite sets $S$ and every integer
twist, which has the corrected proper values, every finite
coefficient comparison, and exactly the rational values E.4g.
It is useful to state the construction and its limits precisely.
Let
\[
 D_X^n=R\operatorname{Hom}
        (R\Gamma(X,\mathbb Q(1-n)),\mathbb Q[-4])
\]
and let $\alpha_X^n:D_X^n\to A_\varnothing^n$ be the
proper correction of \citet[Theorem 3.5]{D_FlachMorin2018}.
For $n\ne1$ define
\[
 C_S^n=
 \operatorname{Cone}\left(
     D_X^n\xrightarrow{\alpha_X^n}A_\varnothing^n
                        \longrightarrow A_S^n\right).
 \tag{E.4.27}
\]
For $n=1$ use E.4o, including its proper member. Choose
functorial cochain representatives once, so all cones and
restriction maps in this definition are actual cochain cones.
For $S=\varnothing$ these are the proper Flach--Morin complexes.
For $n=0$ and $n\ge2$ the correction $D_X^n$ is zero;
thus $C_S^n=A_S^n$ in those weights.

For every $n,S$ and $m\ge1$ there are natural isomorphisms
\[
 C_S^n\otimes^{\mathbb L}\mathbb Z/m
       \simeq A_S^n\otimes^{\mathbb L}\mathbb Z/m
       \simeq
 R\Gamma(\overline U_{S,\mathrm{et}},\mathbb Z/m(n)).
 \tag{E.4.28}
\]
The first isomorphism follows because the correction is
uniquely divisible, and the second because derived reduction
is the cone of multiplication by $m$.
For a single additional removed prime $p$, localization has
the form
\[
 L_p^n\longrightarrow C_S^n\longrightarrow C_{S\cup\{p\}}^n
       \longrightarrow L_p^n[1],
 \qquad p\notin S,
 \tag{E.4.29}
\]
where
\[
 L_p^n\simeq
 \begin{cases}
 C_p^{n-1}[-2]\oplus
   R\Gamma(\mathbb Q_p,\mathbb Q_p/\mathbb Z_p(n))[-2],
       &n<0,\\
 D_p[-3],&n=0,\\
 C_p^0[-2],&n=1,\\
 C_p^{n-1}[-2],&n\ge2.
 \end{cases}
 \tag{E.4.30}
\]
For $n\ne1$ this is the actual etale localization triangle:
the fixed correction $D_X^n$ cancels in the morphism of
defining triangles. For $n=1$ it is E.4e and E.4o.
Multiple removals give the direct sum of these supports and
the compatible iterated localization triangles.

All objects in (E.4.27) are bounded. The proper objects and
their rational correction are bounded by the number-ring
calculation. The new local terms in (E.4.30) are bounded by
the local cohomological dimension bound and E.4a; localization
then proves the assertion for every finite $S$.
Finite generation is deliberately not an assertion:
(E.4.16) proves that it would be false.
Every support in (E.4.30) is rationally zero unless $n=1$.
Consequently (E.4.27) gives the entire rational table E.4g
and its restriction invariance, with the twist-one exception
already calculated in E.4c.

After inverting the residue characteristic $p$, the extra
summands in (E.4.30) vanish and the support becomes
$C_p^{n-1}[-2]\otimes\mathbb Z[1/p]$ in every twist.
In twist zero this follows from the exact unit-reduction
sequence in E.4l, whose quotient is $p$-primary; in negative
twists it follows directly from (E.4.19).
Thus the unrestricted integral purity statement is replaced
by an all-twist statement away from the residue
characteristic, and by the exact wild-support formula at it.

This construction settles existence and the finite-coefficient
comparison for this family of number-ring opens with the
stated coefficients and restrictions. It does not identify
the family with every proposed open Weil-etale topos, nor
construct an additional multiplicative theory on arbitrary
arithmetic schemes. The literal integral values requested in
Lead E2 remain refuted by E.4b and E.4d. The rational
manuscript localization is preserved, and is not the
integral purity assertion refuted in E.4j.
